\uuid{W3ii}
\titre{Ensemble produit}
\auteur{bodin ; hanani ; mzari}
\organisation{exo7}
\licence{CC BY-NC-SA 4.0}
\datecreate{2020-01-20}
\niveau{L1}
\module{Algèbre}
\chapitre{Logique, ensemble, raisonnement}
\sousChapitre{Ensemble}
\difficulte{3}
\theme{Ensembles, applications, ensemble produit}
\exo7qcmid{lille-1.ensembles.021}

\contenu{
\begin{question}
\qclassification{100.02,101.01,102.01,102.02}
\qtags{motcle=ensemble produit}
Soit $A=\{(x,y)\in \Rr^2\mid 2x-y=1\}$ et $B=\{(t+1,2t+1)\mid t\in \Rr\}$. Que peut-on dire de $A$ et $B$ ?
\begin{answers}
    \bad{$A\varsubsetneq B$}
    \bad{$B\varsubsetneq A$}
    \bad{$A\neq B$}
    \good{$A=B$}
\end{answers}
\begin{explanations}
D'abord, $2(t+1)-(2t+1)=1$. Donc $B\subset A$. Réciproquement, pour tout $(x,y)\in A$, il existe $t\in \Rr$ tel que $x=t+1$ et donc $y=2t+1$. D'où $(x,y)\in B$.
\end{explanations}
\end{question}
}
